\begin{textsample}{NEG Dim 4 – 2009\_c5 – Score -1.00 – 2009\_c5/enade\_2009\_com\_soc\_cine\_3a63805496.txt}  \label{ex:f4_neg_082}
Capítulo 5 Distribuição dos Conceitos 5.1 Panorama nacional da distribuição dos conceitos A Tabela 5.1 apresenta a quantidade e o percentual de cursos de COMUNICAÇÃO SOCIAL – CINEMA participantes do ENADE/2009 que se situaram em cada uma das faixas de conceitos, distribuídos de acordo com a Região em que se localizam os cursos das instituições de educação superior, mantenedoras dos cursos responsáveis pela oferta.
O Gráfico 5.1 apresenta a mesma informação sob a forma de diagrama de caixa para uma melhor visualização da dispersão dos conceitos obtidos.
No diagrama de caixa, a caixa propriamente dita fica limitada pelos quartis.
A barra dentro da caixa em negrito representa a mediana.
Linhas verticais unem os quartis aos valores extremos.
Se estes valores extremos se configurarem como valores atípicos, estes aparecem como pontos isolados e as linhas terminam no maior ( ou menor ) valor ainda não considerado atípico.
Eventuais valores atípicos aparecem como pequenos círculos com a Categoria Administrativa correspondente.
Valores considerados muito atípicos foram representados por estrelas.
Um diagrama de caixa para poucos dados ou com dados muito concentrados pode ter alguns pontos coincidentes.
Por exemplo, se mais de metade dos cursos de uma Região receberem conceito 5, o valor extremo superior, o quartil superior e a mediana coincidirão.
Se o conjunto só contém um dado, este apareceria denotado como a barra da mediana.
Se o conjunto apresenta todas as notas iguais, a notação é a mesma, somente uma barra horizontal.
Observando -se os dados da Tabela 5.1, percebe -se que, em termos nacionais, predominaram os cursos que receberam conceito 5, correspondendo, a 40,0 \% do total nacional.
Apenas 1 curso foi classificado com conceito 1 e 1 curso com conceito 2 correspondendo, cada um, a 6,7 \% do total.
O conceito 4 foi \textbf{alocado} a 2 cursos ( 13,3 \% do total nacional ) sendo que 5 ( 33,3 \% ) do total de 15 cursos participantes ficaram sem conceito.
Nenhum curso obteve o conceito 3.
Considerando -se separadamente as cinco Regiões, observou -se que, da Região Norte, nenhum curso foi conceituado.
Na Região Nordeste, 3 cursos participaram, dos quais, nenhum foi avaliado.
Na Região Sudeste, dos 6 cursos participantes, 3 obtiveram o conceito máximo, 2 obtiveram o conceito 4 e 1, o conceito 2.
Na Região Sul, dos 4 cursos participantes, 3 foram avaliados sendo que 2 obtiveram o conceito máximo e 1, o conceito mínimo.
Na Região Centro-Oeste, dos 2 cursos participantes apenas 1 foi avaliado com o conceito máximo. 5.2 Conceitos por Categoria Administrativa e por Região A Tabela 5.2 apresenta a distribuição dos cursos participantes do ENADE/2009 de COMUNICAÇÃO SOCIAL – CINEMA, por Categoria Administrativa, de acordo com os conceitos por eles alcançados, segundo Região.
O Gráfico 5.2 apresenta a mesma informação sob a forma de diagrama de caixa para uma melhor visualização da dispersão dos conceitos obtidos.
Somente eventuais valores atípicos aparecem com o identificador da Região na qual o curso foi habilitado.
De acordo com as informações apresentadas, a maioria, 66,7 \%, é de cursos em instituições privadas, 20,0 \% dos cursos eram vinculados a instituições federais, 13,3 \% a instituições estaduais.
Não houve participação de instituições municipais.
Dos 15 cursos participantes, apenas 10 foram avaliados, ou seja, 5 ficaram sem conceito ( 2 em instituições federais e 3 em instituições privadas ).
O único curso de instituição federal avaliado recebeu conceito 2.
Os 2 cursos de instituições estaduais avaliados, receberam os conceitos mínimo e máximo.
Dos 7 cursos avaliados de instituições privadas, 5 receberam o conceito máximo e 2 o conceito 4.
Não houve participação de cursos na Região Norte.
Na Região Nordeste, dos 3 cursos participantes, 2 eram de instituições federais e 1 de instituição privada e todos ficaram sem conceito.
Na Região Sudeste, dos 6 cursos participantes, 5 eram de instituições privadas e 1 de instituição federal.
Este único curso de instituição federal recebeu o conceito 2 e dos 5 cursos em instituições privadas, 2 receberam o conceito 4 e 3 o conceito máximo.
Na Região Sul, dos 4 cursos participantes, 3 eram de instituições privadas, e 1 de instituição estadual.
Um curso de instituição privada não foi avaliado e os outros 2, receberam conceito máximo.
O único curso de instituição estadual recebeu o conceito mínimo.
Na Região Centro-Oeste, dos 2 cursos participantes, 1 em instituição estadual e o outro em instituição privada, que ficou sem conceito.
O único curso de instituição estadual participante recebeu conceito máximo. 5.3 Conceitos por Organização Acadêmica e por Região Na Tabela 5.3 encontra -se a distribuição dos conceitos atribuídos aos cursos participantes do ENADE/2009 de COMUNICAÇÃO SOCIAL – CINEMA, por Organização Acadêmica, segundo as Regiões.
O Gráfico 5.3 apresenta a mesma informação sob a forma de diagrama de caixa para uma melhor visualização da dispersão dos conceitos obtidos.
Somente eventuais valores atípicos apareceram com o identificador da Categoria Administrativa.
Os cursos desta área estão principalmente \textbf{alocados} em Universidades ( 66,7 \% ) e, Faculdades ( 26,7 \% ) e Centros Universitários ( 6,7 \% ).
Não houve presença de Centros Federais de Educação Tecnológica ( CEFET ) e Institutos Federais de Educação, Ciência e Tecnologia ( IFECT ) e demais categorias.
Nas Universidades, dos 12 cursos participantes, 5 ( 30,0 \% ) não foram avaliados, 4 ( 50,0 \% ) receberam o conceito máximo, 2 ( 10,0 \% ) receberam o conceito 4 ou 1 curso obteve o conceito 2 ( 10,0 \% ).
O único curso de Centro Universitário recebeu conceito 4.
Nas faculdades, dos 7 cursos participantes, 5 ficaram sem conceito, 1 obteve o conceito mínimo e 1 o conceito máximo.
Como já foi mencionado, não foi avaliado nenhum curso desta área na Região Norte.
Na Região Nordeste dos 6 cursos participantes, 2 eram de Universidades e 4 de Faculdade e todos ficaram sem conceito.
Na Região Sudeste, dos 7 cursos participantes, 5 eram de Universidades, 1 de Centro Universitário e 1 de Faculdade.
Nas Universidades, 1 curso obteve o conceito máximo, 2 cursos o conceito 4 e 1 curso o conceito 2.
O curso de Centro Universitário obteve o conceito 4 e o curso em Faculdade ficou com o conceito máximo.
Na Região Sul, dos 4 cursos participantes, 3 eram de Universidades e 1 de Faculdade.
Nas Universidades, 2 cursos obtiveram o conceito máximo e 1 curso não foi avaliado.
O curso de Faculdade obteve o conceito mínimo.
Na Região Centro-Oeste, dos 3 cursos participantes, 2 eram de Universidades e 1 de Faculdade.
Um curso de Universidade obteve o conceito máximo, o outro ficou sem conceito e o curso de Faculdade não foi avaliado.

% matched lemmas: alocar
\end{textsample}
